\documentclass[11pt]{article}
\usepackage[margin=1in]{geometry}
\usepackage{amsmath,amssymb,amsthm}
\usepackage[hidelinks]{hyperref}
\newtheorem{theorem}{Theorem}
\newtheorem{lemma}[theorem]{Lemma}
\newcommand{\HM}{\mathscr H}
\usepackage{mathrsfs}
\title{Discontinuity of the weighted harmonic Hardy constant\\
\large A counterexample to MathDB problem 345844}
\author{Principia Math in Progress}
\date{September 12, 2026}
\begin{document}
\maketitle

\paragraph{Task 1. Resolve the continuity conjecture.}
Conjecture 1 of Pasteczka \cite[p.~1061]{P} asserts that a symmetric,
monotone weighted mean, continuous in its weights and having the
Hardy property, has a Hardy constant continuous for pointwise
convergence of positive weight sequences. We refute this with the
weighted harmonic mean on \(I=(0,\infty)\).

For positive weights \(w=(w_n)_{n\ge1}\), put \(W_n=\sum_{i=1}^n w_i\) and
\[
 M_n(x;w)=\frac{W_n}{\sum_{i=1}^n w_i/x_i}.
\]
Write \(\HM(w)\) for the least \(C\in[0,\infty]\) such that
\[
 \sum_{n\ge1}w_nM_n(x;w)\le C\sum_{n\ge1}w_nx_n
 \quad\text{whenever }x_n>0,\quad \sum_{n\ge1}w_nx_n<\infty.
 \tag{1}
\]
This is the source's Definition 2.3. The harmonic mean is symmetric
under simultaneous permutation of inputs and weights, monotone in
each input, and continuous in the nonnegative weights when their
sum is positive. Its formula also gives weight homogeneity, reduction,
zero-weight elimination, and the minimum--maximum bounds required
of a weighted mean.

\begin{lemma}\label{lem:upper}
For every positive weight sequence \(w\), one has \(\HM(w)\le9/2\).
\end{lemma}
\begin{proof}
The Cauchy--Schwarz inequality and a right-endpoint integral bound give
\[
 \left(\sum_{i=1}^n w_i\sqrt{W_i}\right)^2
 \le\left(\sum_{i=1}^n\frac{w_i}{x_i}\right)
       \left(\sum_{i=1}^n w_iW_ix_i\right),\qquad
 \sum_{i=1}^n w_i\sqrt{W_i}\ge\int_0^{W_n}\sqrt t\,dt
 =\frac23W_n^{3/2}.
\]
Consequently,
\[
 M_n(x;w)\le\frac{9}{4W_n^2}\sum_{i=1}^n w_iW_ix_i.
 \tag{2}
\]
Since \(t^{-2}\) decreases, for every \(i\ge1\),
\[
 \sum_{n=i}^{\infty}\frac{w_n}{W_n^2}
 \le\frac{w_i}{W_i^2}
   +\sum_{n=i+1}^{\infty}\int_{W_{n-1}}^{W_n}\frac{dt}{t^2}
 \le\frac2{W_i}.
 \tag{3}
\]
Tonelli's theorem, (2), and (3) show that the left side of (1) is at most
\[
 \frac94\sum_{i\ge1}w_iW_ix_i
             \sum_{n\ge i}\frac{w_n}{W_n^2}
 \le\frac92\sum_{i\ge1}w_ix_i.
\]
In particular the harmonic mean satisfies the conjecture's Hardy
hypothesis, even in the formulation requiring every positive
weight sequence to have a finite Hardy constant.
\end{proof}

\begin{lemma}\label{lem:lower}
If \(0<w_n\le B<\infty\) and \(W_n\to\infty\), then \(\HM(w)\ge2\).
\end{lemma}
\begin{proof}
Put \(Q_n=\sum_{i=1}^n w_i^2\). The identity
\[
 \sum_{i=1}^n w_iW_i=\frac{W_n^2+Q_n}{2}
 \tag{4}
\]
follows by separating diagonal and off-diagonal terms in \(W_n^2\).
Fix \(L\ge1\) and \(\delta>0\), and use the strictly positive input
\[
 x_i=\frac1{W_i}\ (i\le L),\qquad
 x_i=\frac{\delta\,2^{-(i-L)}}{w_i}\ (i>L).
\]
Its weighted sum is \(D_L+\delta\), where \(D_L=\sum_{n=1}^Lw_n/W_n\).
By (4), its first \(L\) output terms have weighted sum
\[
 T_L=\sum_{n=1}^L\frac{2w_nW_n}{W_n^2+Q_n}.
\]
Thus \(\HM(w)\ge T_L/(D_L+\delta)\), and letting \(\delta\downarrow0\)
gives \(\HM(w)\ge T_L/D_L\).

For \(t\in[W_{n-1},W_n]\) we have \(t+B\ge W_n\). Hence
\[
 D_L\ge\int_0^{W_L}\frac{dt}{t+B}
      =\log(1+W_L/B)\longrightarrow\infty.
\]
On the other hand \(Q_n\le BW_n\), so (3) with \(i=1\) yields
\[
 0\le2D_L-T_L
 =\sum_{n=1}^L\frac{2w_nQ_n}{W_n(W_n^2+Q_n)}
 \le2B\sum_{n=1}^L\frac{w_n}{W_n^2}
 \le\frac{4B}{w_1}.
\]
It follows that \(T_L/D_L\to2\), proving the assertion.
\end{proof}

\begin{theorem}
The conjectured continuity fails, even among strictly positive
integer weight sequences with infinite total mass.
\end{theorem}
\begin{proof}
For \(n,K\ge1\), define
\[
 \lambda_n=4^{\,n-1},\qquad
 \mu_n^{(K)}=
 \begin{cases}
 4^{\,n-1},&n\le K,\\
 4^{\,K-1},&n>K.
 \end{cases}
\]
Both sequences have infinite total mass, and
\(\mu^{(K)}\to\lambda\) pointwise: each fixed coordinate agrees
exactly once \(K\ge n\). Each \(\mu^{(K)}\) is bounded, so
Lemma \ref{lem:lower} gives \(\HM(\mu^{(K)})\ge2\).
For the limiting sequence,
\[
 M_n(x;\lambda)
 \le\frac{\Lambda_n}{\lambda_n}x_n
 =\frac{4-4^{\,1-n}}3x_n\le\frac43x_n.
\]
Multiplication by \(\lambda_n\) and summation give
\(\HM(\lambda)\le4/3\). Thus
\[
 \HM(\mu^{(K)})-\HM(\lambda)\ge\frac23
 \quad\text{for every }K,
\]
which contradicts continuity. Lemma \ref{lem:upper} and the stated
properties of the harmonic mean verify every hypothesis.
\end{proof}

\paragraph{Sharper values and verification.}
The proof above is self-contained and needs no sharp Hardy theorem.
As a consistency check, the source's Lemma 4.1, with Gini parameters
\((p,q)=(-1,0)\), gives \(\HM(\lambda)=5/4\) and
\(\HM(\mu^{(K)})=2\): the respective nonincreasing ratios
\(w_n/W_n\) tend to \(3/4\) and \(0\).
The accompanying exact rational checker verifies finite positive-input
witnesses, identity (4), and the stated bounds. Finite tests are
supplementary; the proof covers every \(K\) and infinite sequences.

\paragraph{Task 2. Classify weights attaining the maximal constant.}
Find explicit necessary and sufficient conditions on a positive
weight sequence for its harmonic Hardy constant to equal \(2\),
including irregular sequences outside the monotone-ratio setting.

\paragraph{Task 3. Determine constants for oscillating weight growth.}
Compute the harmonic Hardy constant when \(w_n/W_n\) has two or
more accumulation points, and determine which features beyond those
points affect the answer.

\paragraph{Task 4. Find useful topologies restoring continuity.}
Identify a topology stronger than coordinatewise convergence that
permits meaningful infinite-tail perturbations and makes weighted
Hardy constants continuous for broad classes of Hardy means.

\paragraph{Task 5. Quantify the effect of a tail replacement.}
For harmonic and other power means, bound the change in the Hardy
constant under controlled perturbations of all weights after a
specified index, with explicit dependence on the perturbation size.

\begin{thebibliography}{9}
\bibitem{P}
P.~Pasteczka, \emph{On properties of weighted Hardy constant for means},
Mathematical Inequalities \& Applications \textbf{25}(4) (2022),
1047--1068.
\href{https://doi.org/10.7153/mia-2022-25-65}{doi:10.7153/mia-2022-25-65};
\href{https://files.ele-math.com/articles/mia-25-65.pdf}{published text};
\href{https://arxiv.org/abs/2003.06025}{arXiv:2003.06025}.
Definitions: pp.~1051--1052. Conjecture 1 and Lemma 4.1: p.~1061.
\end{thebibliography}
\end{document}
